% !PRINTED 18
\weberpage{18}
%<<CONT
\begingroup
\setlength{\abovedisplayskip}{5pt plus 2pt minus 1pt}
\setlength{\belowdisplayskip}{5pt plus 2pt minus 1pt}
\setlength{\abovedisplayshortskip}{3pt plus 1pt}
\setlength{\belowdisplayshortskip}{4pt plus 1pt}
and we further stipulate that $(x,0)=x$, which does not contradict
these equations. The pair $(x,y)$ is $=0$ only when $x$ and $y$
are both zero. For abbreviation we denote $(0,1)$ by $i$. The
equations above then give the consequences:
\[
 (x,0)(0,1)=(0,x)\quad\text{or}\quad =ix,
\]
\[
 (x,0)+(0,y)=(x,y)=x+yi,
\]
\[
 x+yi+a+bi=(x+a)+(y+b)i,
\]
and the inverse of addition:
\[
 x+yi-(a+bi)=(x-a)+i(y-b),
\]
and further multiplication:
\[
 (x+yi)(a+bi)=xa-yb+i(xb+ya),
 \qquad i^2=-1,
\]
\[
 (x+yi)(x-yi)=x^2+y^2,
\]
\[
 x+yi=(a+bi){(a-bi)(x+yi)\over a^2+b^2}
\]
or
\[
 {x+yi\over a+bi}={ax+by+i(ay-bx)\over a^2+b^2},
\]
which defines division except when $a+bi=0$.

Numbers of the form $ix$ are called \gesperrt{purely imaginary
numbers}, and $i$ the imaginary unit. The numbers $a+bi$ are called
\gesperrt{imaginary} or \gesperrt{complex}. The systems of
\gesperrt{real} and \gesperrt{purely imaginary numbers} are contained
among them as special cases.

Thus, in the domain of \gesperrt{complex numbers} $x+yi$, the
fundamental operations can be carried out without restriction (except
for division by zero), and calculation with real numbers is a special
case.

Following \gesperrt{Gauss}, one represents the complex numbers
$z=x+yi$ geometrically by the points of a plane, laying down a
rectangular coordinate system and regarding the point with coordinates
$x,y$ as the image of the numerical value $z$. The points of the
$x$-axis represent the real numbers $x$ in the way discussed above.
The points of the $y$-axis are the images of the purely imaginary
numbers $yi$. The origin is the image of the number $0$. The radius
vector from the origin to the point $z$ has the numerical value
$\varrho=\sqrt{x^2+y^2}$, which is called the \gesperrt{absolute
value}, or the magnitude, or, according to older terminology, the
modulus of the complex number $z$.
\endgroup
